\chapter{Related Work}
\label{chapter:related_work}

Research on modeling population dynamics spans two central but historically distinct areas: crowd simulation and population mobility modeling. Crowd simulation emphasizes fine-grained representations of movement and interactions among individuals or groups, whereas population mobility focuses on larger-scale flows and statistical regularities in urban dynamics. Each domain has produced significant advances, but both exhibit limitations that motivated the development of LodusPop.

\section{Crowd Simulation}
\label{related:crowd_simulation}

Crowd simulation has long served as a foundation for understanding human movement in both indoor and outdoor environments, with models developed at microscopic, macroscopic, and hybrid scales.

Microscopic models treat every agent as an autonomous entity with decision-making capabilities, typically influenced by local environmental conditions. Classical examples include the Social Force Model~\cite{helbing1995social}, where pedestrian motion emerges from forces balancing personal goals with environmental and social interactions, and ORCA~\cite{van2011reciprocal}, which guarantees collision-free trajectories through reciprocal velocity obstacles. BioCrowds~\cite{de2012simulating} introduced a space-discretization strategy where agents compete for territory, naturally generating emergent Voronoi-like partitions. Later studies aimed to add heterogeneity by embedding psychological or emotional states into agent behavior. For example, Pelechano et al.~\cite{pelechano2005crowd} and Durupinar et al.~\cite{durupinar2015psychological} integrated factors such as stress, emotion, and cultural variation, while Neto et al.~\cite{neto2017giving} modeled emotional contagion based on personality traits. These approaches achieve remarkable behavioral variety, but they depend heavily on detailed behavioral datasets and quickly become computationally expensive when scaled up to large urban populations.

In contrast, macroscopic models treat crowds as collective entities, often described through density fields or flow equations. Hughes~\cite{hughes2002continuum} and Treuille et al.~\cite{treuille2006continuum} pioneered continuum models where population density evolves through potential fields defined over the environment, while proposed Antonitsch et al.~\cite{antonitsch2019bioclouds} BioClouds,
aggregating agents into homogeneous “clouds” competing for space based on density and velocity preferences. Such formulations shift
part of the steering to environment-side representations and significantly improve computational scalability, but they sacrifice the diversity of individual behaviors.

Hybrid models attempt to combine these two perspectives. For example, Legion\cite{legion} dynamically adapts the level of detail based on environmental complexity, allowing legions of agents to split or merge as situations demand. Xiong et al.~\cite{xiong2010hybrid} proposed spatially partitioned hybrid simulations that combine ORCA~\cite{van2011reciprocal} in high-interest regions with continuum models elsewhere, while Lemonari et al.~\cite{lemonari2025mpact} introduced MPACT, a mesoscopic framework that uses learned profiles to capture intermediate-scale behaviors. These hybrids achieve flexibility, however, the primary focus of
control remains on the population side. Agents query these fields and detailed data is often required for accurate calibration.

Recent advances explore ways to alleviate this dependency through data-driven modeling, where behavior is learned from observed or simulated trajectories rather than being entirely hand-crafted. For instance, machine learning has been used to embed agent dynamics into low-dimensional manifolds, enabling multiscale reconstruction of crowd flows~\cite{alvarez2025quationfree}. 
Deep reinforcement learning approaches, such as those proposed by Li et al.~\cite{Li2025efficienteCrowdSim}, combine learned navigation strategies with anisotropic environmental fields to reduce computational costs in large spaces. 
In many of these studies, policies are learned from large numbers of simulated episodes rather than from real mobility datasets. Behavior is governed by fitted models whose data requirements center on simulated experience rather than fine-grained observational traces, and the underlying agent-based perspective remains unchanged.

In summary, crowd simulation excels in behavioral realism and the reproduction of fine-scale interactions. However, the reliance on detailed micro-data and high computational costs make it difficult to deploy at the scale of entire cities or to apply when granular behavioral data is unavailable.


\begin{comment}
Crowd simulation techniques are commonly used to represent the movement of a population, both in interior and exterior environments and in various levels of detail, including microscopic points of view~\cite{helbing1995social, pelechano2005crowd, narang2015generating, cao2021competitiveCooperative, zhang2023collectiveBehavior}; macroscopic points of view~\cite{hughes2002continuum, treuille2006continuum}; and a combination of both (i.e., hybrid)~\cite{antonitsch2019bioclouds}. 

% Microscopic models with homogeneous behaviors
In microscopic crowd simulation models, each agent in a crowd is a single entity of simulation\cite{macal2010agentBasedTutorial, xiong2014motionPlanning}. Some models propose homogeneous behaviors ~\cite{helbing1995social, van2011reciprocal, de2012simulating}, where agents will behave similarly to each other and show little to no individuality. 
Helbing and Molnár~\cite{helbing1995social} presented the Social Forces Model, where agent movement is influenced by their desired velocity toward a destination, distance from other agents, and interaction with the environments (e.g., obstacles and walls). 
ORCA (Optimal Reciprocal Collision Avoidance), proposed by Van den Berg et al.~\cite{van2011reciprocal}, simulates agents based on reciprocal velocity obstacles (RVO).
BioCrowds, presented by Bicho et al.~\cite{de2012simulating}, proposes a space discretization and competition model for crowd simulation, where agents compete to maintain their personal space by taking possession of space around them, forming an emergent Voronoi partitioning of space. A collision-free movement is guaranteed if no agent attempts to move outside its individual space.

% Microscopic models with heterogeneous behaviors (could be removed/reduced/combined with previous paragraph)
There has been research on how to model heterogeneous behavior for individual agents in crowds.
Pelechano et al.~\cite{pelechano2005crowd} proposed a model to incorporate psychological aspects and communication into crowd simulations to generate emergent cultural behavior. The authors model human behavior like stress, emotion, motivation, decision process, and ability, impacting decision-making under the agent's varied emotional state (e.g., stress levels and physiological state).
Durupinar et al.~\cite{durupinar2015psychological} proposed a model to differentiate crowds between audiences and mobs as passive and active crowds, respectively, influenced by emotional and homogeneous behavior. Authors present an emotion contagion model based on OCEAN personality traits (\textit{Openness}, \textit{Conscientiousness}, \textit{Extroversion}, \textit{Agreeableness} and \textit{Neuroticism}), employing the OCC emotion model in multi-agent simulations, and utilizing a pleasure-arousal-dominance (PDA) model do influence agent decision-making.
Neto et al.~\cite{neto2017giving} propose a model based on BioCrowds~\cite{de2012simulating} to simulate the spread of psychological phenomena in a crowd of virtual agents. Agents are imbued with OCEAN personality traits and emotional information and can spread that information based on factors such as expressiveness, susceptibility to contagion, and current goal.


% Macroscopic models
In macroscopic crowd simulation models, agents are represented as part of a collective being controlled or influenced by a macro-level structure~\cite{thalmann2007crowd}, e.g., density fields or aggregations of agents. This allows a higher granularity and smaller simulation cost with a trade-off of losing individual characteristics and behaviors for each agent. 
Treuille et al.~\cite{treuille2006continuum} proposed the Continuum Crowds model, which is based on continuum field dynamics and inspired by the work of Hughes~\cite{hughes2002continuum}. The authors simulate a crowd as a dynamic potential field to perform navigation for agents and moving obstacles without explicit agent-based dynamics. 
Antonitsch et al.~\cite{antonitsch2019bioclouds} presented BioClouds, based on BioCrowds~\cite{de2012simulating}, which models crowds as a set of clouds. Each cloud comprises a group of agents with a similar desire for a specific density, speed, and goal. Clouds compete for space amongst each other and occupy the environment in a manner that tries to keep their desired densities respected. 

% Hybrid models
Hybrid approaches for crowd simulation have also been proposed, combining distinct microscopic and macroscopic models and leveraging the accuracy and computational efficiency. Xiong et al.~\cite{xiong2010hybrid} suggest separating the environment into static regions. Regions of interest are simulated with ORCA~\cite{van2011reciprocal}, and the remaining space with Continuum Crowds~\cite{treuille2006continuum}. The regions are interfaced with transition regions, with agents that cross these boundaries being removed from one simulation and added to the other. 
BioClouds~\cite{antonitsch2019bioclouds} also presents a way to transition between macroscopic and microscopic simulations. The authors utilize a visualization window that covers a region in the environment. The area within this region and agents of clouds that occupy it are simulated with the BioCrowds~\cite{de2012simulating} model.
Legion, proposed by Antonitsch et al.\cite{legion}, presents a crowd simulation model where the simulation dynamically changes the level of detail of the modeled agents according to environment complexity. A legion (i.e., crowds of people) may split into smaller groups whenever it's in a situation where it would occupy less space than desired in a future frame, e.g., moving towards an obstacle. Different legions may also merge when they share the same objective and are close enough to regroup.
\red{
Lemonari et al.~\cite{lemonari2025mpact} propose MPACT, a mesoscopic crowd simulation framework that transforms unlabelled trajectory data into interpretable behavior profiles using image-based encodings. Unlike traditional microscopic or macroscopic models, MPACT operates at an intermediate level, capturing group dynamics and localized interactions. Trained on synthetic data, its predictive model maps spatial encodings to high-level simulator parameters, enabling behavior-driven simulations without manual tuning.}
\end{comment}

\section{Population Mobility}

Population mobility studies emerged from statistical and large-scale analyses of human movement patterns. Unlike crowd simulation, which centers on the agent, mobility modeling typically abstracts away individuality to emphasize flows across geographic space and time.

A consistent finding in this field is that human displacements follow heavy-tailed distributions. Brockmann et al.~\cite{brockmann2006scaling} showed that the circulation of banknotes in the United States resembled Lévy flights, while González et al.~\cite{gonzalez2008understanding}, analyzing cell phone data, demonstrated that most individuals repeatedly visit a small set of locations. Song et al. ~\cite{natureSong2010exploration}  formalized this as a combination of exploration (discovering new places) and preferential return (frequent revisits), and Alessandretti et al.~\cite{alessandretti2020nature} refined these ideas with hierarchical ``containers'', arguing that people conceptualize space in nested levels such as neighborhoods, cities, and regions. These statistical regularities provide strong predictive power, but they are generally limited to describing long-term patterns rather than simulating short-term interventions or local policies.

The COVID-19 pandemic highlighted the importance of connecting mobility with public health. Studies such as Jamshidi et al.~\cite{jamshidi2020global} examined the effect of density and mobility restrictions on infection rates, while Hernández-Orallo et al.~\cite{hernandez2021human} used urban simulations to identify risk hotspots like train stations and plazas. Kraemer et al.~\cite{natureKramer2020mobilityCovid} quantified how mobility and control measures influenced epidemic spread in China. These works demonstrated the societal relevance of mobility modeling but remained largely domain-specific, oriented toward epidemiology rather than more general urban planning applications.

Frameworks for simulating mobility have also emerged. Silva et al.~\cite{silva2020lodus} introduced a framework that integrates environmental data with macroscopic BioClouds to simulate contagion scenarios. Although innovative, its reliance on homogeneous clouds and a narrow focus on epidemics limited its applicability. More recent frameworks increasingly incorporate machine learning and multi-scale techniques [6,7,16], but they remain strongly tied to data availability and typically adhere to the population-driven logic of traditional mobility models.

Recent contributions push the field toward more interpretable, adaptable, and AI-enhanced approaches. Liu et al.~\cite{liu2024GraviGBM} proposed GraviGBM, which blends gravity-based migration models with gradient-boosted machine learning, balancing predictive accuracy and interpretability. Other work has used symbolic regression to uncover interpretable mobility laws directly from data, identifying power-law patterns in flows~\cite{guo2025datadrivendiscoveryhuman}. %Cabanas-Tirapu et al. [2025] showed that closed-form machine-learned models can match the predictive accuracy of deep learning while remaining interpretable. In parallel, human mobility has become a key ingredient in calls for spatial foundation models that move beyond static points of interest to dynamic, place-based representations [Hashemi \& Züfle, 2025]. 
On the simulation side, new agent-based frameworks integrate large language models with mobility environments, as in CAMS~\cite{du2025camscitygptpoweredagenticframework}, enabling autonomous agents to generate plausible movement patterns without pre-defined geospatial inputs.

Additionally, several general-purpose urban simulation frameworks are widely used. MATSim~\cite{balmer2009matsim} is a large-scale agent-based framework for transportation planning that simulates individual travel demand and traffic assignment through iterative optimization. SUMO~\cite{krajzewicz2012sumo} is a microscopic traffic simulator designed for modeling vehicles and road networks, offering fine-grained mobility but limited population abstraction. GAMA~\cite{taillandier2019gama} provides a flexible multi-agent platform where users can model a wide range of urban and environmental processes with rich GIS integration.
%While these frameworks are powerful, they rely on population-driven dynamics, requiring either detailed mobility traces (MATSim) or pre-specified decision rules (GAMA), making them heavily dependent on granular data. Also, they are often domain-specific: SUMO is tailored to traffic, MATSim to transportation demand, and GAMA to general-purpose agent-based modeling. %By contrast, LodusPop introduces an environment-driven paradigm, where environments request population flows, reducing the reliance on fine-grained mobility data. Its abstraction of populations into dynamic “blobs” allows simulation across multiple spatial scales, something traditional frameworks only achieve through labor-intensive model specification.

In summary, mobility research provides robust statistical frameworks and promising AI-driven methods, but most approaches are either too aggregated for dynamic, multi-scale interventions or remain dependent on detailed mobility data.


\begin{comment}
Population mobility presents an abstraction of key aspects of a population from a set of parameters, where the notion of the individual is often non-existent, and the information available is related to groups of people. This abstraction is usually a statistical approach to large-scale temporal human mobility and migration patterns, with analytical studies presented on the characteristics of human mobility in urban environments~\cite{brockmann2006scaling,gonzalez2008understanding}. 

Due to the recent COVID-19 pandemic, there has been a renewed interest in the study of urban mobility and its effects on disease spreading~\cite{jamshidi2020global,hernandez2021human, natureKramer2020mobilityCovid}.
Brockmann et al.~\cite{brockmann2006scaling} proposes that human mobility follows a power law distribution. The authors tracked the displacement of over $400$ thousand banknote serial numbers with over a million reports of their locations throughout the United States. The trajectories of these bank notes are reminiscent of Lévy flights, a scale-free random walk\cite{levy_flight,chechkin2008levy}.
Gonzalez et al.~\cite{gonzalez2008understanding} also proposed that mobility can be presented as a power law. The authors used data from cell phones, where their positions are defined within the area of the closest reception tower when a phone call starts. A displacement is detected only when a phone is detected within the area of a different tower. This data is a better representation of the movement of individuals, as cell phones are a more personal item than bank notes, which are often exchanged between people.
Indeed, many studies found in the literature utilize power laws (followed in frequency by exponential and log-normal distributions) to describe human mobility, as presented by Alessandreti et al.~\cite{natureAlessandretti2017multiScaleSpatioTemporal}. This study also proposed that the waiting times between the displacement of people follow a log-normal distribution.

One of the key findings of these studies~\cite{brockmann2006scaling, gonzalez2008understanding} is that humans mostly visit a small set of locations repeatedly in their daily routine, with a random chance of visiting novel locations as time passes. Additionally, the authors conclude that people in a city tend to travel more frequently to places closer to their current location, with fewer longer jumps.
Song et al.~\cite{natureSong2010exploration} proposed a model to understand and predict the mobility of individuals based on some of these findings. The authors argue that, instead of simply following a continuous time random walk (CTRW) or Lévy flight, people also present behaviors of exploration, where they tend to explore new locations over time (e.g., near their homes or workplaces) and preferential return, where they often return to frequently visited sites. 

In a more recent study, Alessandretti et al.~\cite{alessandretti2020nature} discuss the seeming contradiction between our common mental representation of spaces as a hierarchical structure~\cite{naturePaasi2004placeAndRegion} and the recent findings that population mobility follows scale-free distributions. People often describe spatial entities in association with a given scale or size (e.g., buildings, neighborhoods, cities) and their relations to other spatial entities in a nested structure~\cite{natureBatty2006hierarchy, naturePumain2006hierarchy}. 
The authors propose that human mobility contains meaningful scales, modeled as spatial containers representing similar spatial entities as described before, forming a hierarchy. The observed scale-free distributions emerge from the aggregation of displacements across these containers. The study also indicates that the containers model presents a better fit to mobility data when compared to the exploration and preferential return model of Song et al.~\cite{natureSong2010exploration}.

Regarding disease spreading, Jamshidi et al.~\cite{jamshidi2020global} present an evaluation of the effect of the weather on human mobility, urban density, and social isolation and the subsequent impact on the spread of COVID-19. The work highlights the effects of spatial and temporal scales on the infection models for COVID-19. 
Hernández-Orallo et al.~\cite{hernandez2021human} explored using urban simulation to evaluate the temporal and spatial risks of infection in an urban scenario. The analyzed situations included train stations and public squares, and the authors generated heatmaps to locate areas of increased risk. The goal is to help health agencies mitigate risks in public spaces possibly.

\red{Silva et al.~\cite{silva2020lodus} introduced LODUS, a simulation framework that combines environmental data with macroscopic agent behaviors to model large-scale population dynamics. Their approach utilizes BioClouds\cite{antonitsch2019bioclouds} to simulate the population. The study emphasizes contagion scenarios, explicitly modeling agent routines and social interactions (using BioCrowds~\cite{de2012simulating}). Although the framework provides detailed environmental modeling, its application is primarily demonstrated in the context of contagion using homogeneous BioClouds clouds, leaving its applicability to other domains less explored.}
\end{comment}

\section{Positioning LodusPop Between Crowd Simulation and Mobility Models}

The works reviewed above point to a gap between crowd simulation and population mobility models. On one side, crowd simulators are capable of reproducing rich local interactions and complex behaviors, but they become expensive at city scale and often presuppose detailed, behavior-level data for calibration. Mobility models, in turn, describe flows at coarser levels and scale more easily to large regions, but offer limited support for representing how specific locations, policies, or events structure movement over time in a way that remains explicit and manipulable.

LodusPop is designed to bridge this gap by making the environment an explicit driver of population dynamics while retaining heterogeneity in the population. Routines are attached to regions and points of interest instead of individuals, so locations express which types of population they demand or release at different times. For example, schools requesting children from nearby residential areas, hospitals attracting patients with particular attributes, or shelters drawing specific profiles subject to capacity constraints. 
Populations are represented as blobs of variable size that carry typed, traceable attributes and can split when only a subpopulation participates in an action or undergoes a change of state, and merge again when compatible groups share the same location. In this way, a single simulation can move smoothly between aggregate flows and more detailed descriptions where fine tracking is important, without losing consistency of attributes over time.

The dynamics of LodusPop are expressed in a compact yet extensible action formalism that covers both movements in the environment and changes in traceable characteristics. 
Routines attached to locations are written as sequences of such actions, which the simulator interprets to update blobs and their distribution across the environment. 
Additional processes such as disease transmission, recovery, vaccination, or participation in specific events can be introduced as new action types together with their consumption rules, implemented as plugins that integrate with the same core abstraction. 
As a result, LodusPop can combine different mobility and process models within a unified environment-driven representation, reusing the same engine to explore diverse scenarios while maintaining a clear separation between the core simulator and domain-specific mechanisms.


% Antiga 2.3
\begin{comment}
    {
The reviewed works reveal a consistent gap. Crowd simulation approaches achieve behavioral variety but struggle with scalability and require extensive behavioral data, while population mobility models provide large-scale regularities but lack adaptability for dynamic interventions. Both paradigms overwhelmingly adopt population-driven logic, where individuals determine their own movements. \redb{Even in formulations that encode steering as vector, potential, or interaction fields over the environment~\cite{hughes2002continuum, treuille2006continuum, antonitsch2019bioclouds}, agents still query these fields locally and decide how to move in response.}

LodusPop departs from this paradigm through an environment-driven approach: environments request populations with specific attributes, reducing reliance on granular behavioral data and enabling simulations even where mobility traces are unavailable. By representing populations as flexible ``blobs'' that can split, merge, and evolve, LodusPop fluidly adapts across multiple spatial resolutions, from neighborhoods to entire cities.

Compared to established frameworks such as SUMO, MATSim, and GAMA, LodusPop introduces three distinguishing features. First, unlike SUMO or MATSim, which require detailed travel demand or traffic data, LodusPop can operate with heterogeneous and incomplete data sources such as census statistics or synthetic populations. Second, unlike generic agent-based platforms such as GAMA, where modelers must explicitly encode behavioral rules, LodusPop embeds an environment-driven paradigm that shifts the burden of decision-making from individuals to the environment itself. Finally, the blob abstraction provides a balance between computational efficiency and behavioral expressiveness, allowing population dynamics to be simulated at scales that are impractical for microscopic models.

In this way, LodusPop positions itself as a complementary tool to existing urban simulation frameworks: less data-hungry than agent-based traffic simulators, more adaptable than statistical mobility models, and uniquely suited for rapid policy evaluation across diverse urban contexts.
}
\end{comment}